\section*{第 2 章：数学预备知识}
\addcontentsline{toc}{section}{第 2 章：数学预备知识}

\begin{taplsolutionentrybox}
\phantomsection
\label{sol:translator:2.2.6}
\textbf{练习 2.2.6}

要证明 \(R'\) 是 \(R\) 的自反闭包，需要核对三个条件：它包含 \(R\)，
它是自反关系，并且它是满足前两个条件的最小关系。

\begin{enumerate}[label=\textup{(\arabic*)},leftmargin=2.2em]
  \item 由
  \[
    R'=R\cup\{(s,s)\mid s\in S\}
  \]
  立即得到 \(R\subseteq R'\)。

  \item 任取 \(s\in S\)。根据右侧集合的定义，
  \((s,s)\in\{(s,s)\mid s\in S\}\)，因而 \((s,s)\in R'\)。
  所以 \(R'\) 是自反关系。

  \item 设 \(R''\) 是 \(S\) 上任意一个包含 \(R\) 的自反关系。
  由 \(R\subseteq R''\)，可知 \(R''\) 包含 \(R\) 中的所有有序对；
  又因为 \(R''\) 自反，所以对每个 \(s\in S\) 都有
  \((s,s)\in R''\)。因此
  \[
    R\cup\{(s,s)\mid s\in S\}\subseteq R'',
  \]
  即 \(R'\subseteq R''\)。
\end{enumerate}

由此，\(R'\) 是包含 \(R\) 的最小自反关系，所以它正是 \(R\) 的
自反闭包。

\taplsolutionbacklink{ex:2.2.6}{返回练习 2.2.6}
\end{taplsolutionentrybox}

\begin{taplsolutionentrybox}
\phantomsection
\label{sol:translator:2.2.7}
\textbf{练习 2.2.7}

先注意 \(R_i\subseteq R_{i+1}\)，所以序列
\(R_0,R_1,\ldots\) 单调递增。

\begin{enumerate}[label=\textup{(\arabic*)},leftmargin=2.2em]
  \item 因为 \(R=R_0\subseteq\bigcup_iR_i=R^+\)，所以 \(R^+\)
    包含 \(R\)。

  \item 证明 \(R^+\) 具有传递性。若
    \((s,t)\in R^+\) 且 \((t,u)\in R^+\)，则存在 \(i,j\)，使
    \((s,t)\in R_i\) 且 \((t,u)\in R_j\)。令 \(k=\max(i,j)\)。
    由单调性，两对都属于 \(R_k\)，所以按 \(R_{k+1}\) 的定义，
    \((s,u)\in R_{k+1}\subseteq R^+\)。

  \item 证明最小性。这里的“最小”按关系之间的包含关系理解。设 \(Q\)
    是任意一个包含 \(R\) 的传递关系；我们需要证明
    \[
      R^+\subseteq Q
    \]
    因为 \(R^+=\bigcup_iR_i\)，只要对 \(i\) 作归纳，证明每一层都满足
    \(R_i\subseteq Q\) 即可。

    \textbf{基例。}由 \(R_0=R\) 以及假设 \(R\subseteq Q\)，立即得到
    \(R_0\subseteq Q\)。

    \textbf{归纳步骤。}假设 \(R_i\subseteq Q\)，下面证明
    \(R_{i+1}\subseteq Q\)。任取 \((s,u)\in R_{i+1}\)。根据
    \(R_{i+1}\) 的定义，分为两种情形：
    \begin{enumerate}[label=\textup{\alph*.},leftmargin=2.2em]
      \item 若 \((s,u)\in R_i\)，则由归纳假设直接得到
        \((s,u)\in Q\)。

      \item 否则，存在某个 \(t\)，使
        \((s,t)\in R_i\) 且 \((t,u)\in R_i\)。由归纳假设，
        \((s,t),(t,u)\in Q\)；再由 \(Q\) 的传递性，得到
        \((s,u)\in Q\)。
    \end{enumerate}
    两种情形下都有 \((s,u)\in Q\)，所以
    \(R_{i+1}\subseteq Q\)。归纳可知，对所有 \(i\) 都有
    \(R_i\subseteq Q\)，因而
    \[
      R^+=\bigcup_iR_i\subseteq Q
    \]
    由于 \(Q\) 是任意一个包含 \(R\) 的传递关系，\(R^+\) 包含于所有
    这样的关系，所以它是其中最小的一个。
\end{enumerate}

所以 \(R^+\) 是包含 \(R\) 的最小传递关系，即 \(R\) 的传递闭包。

\taplsolutionbacklink{ex:2.2.7}{返回练习 2.2.7}
\end{taplsolutionentrybox}

\begin{taplsolutionentrybox}
\phantomsection
\label{sol:translator:2.2.8}
\textbf{练习 2.2.8}

定义 \(S\) 上的关系
\[
  Q=\{(s,t)\in S\times S\mid P(s)\Rightarrow P(t)\}
\]
关系 \(Q\) 是自反的，因为 \(P(s)\Rightarrow P(s)\)；它也是传递的，因为
\[
  P(r)\Rightarrow P(s),\quad P(s)\Rightarrow P(t)
  \quad\Longrightarrow\quad
  P(r)\Rightarrow P(t)
\]
“\(P\) 被 \(R\) 保持”恰好说明 \(R\subseteq Q\)。而 \(R^*\) 是包含
\(R\) 的最小自反传递关系，所以 \(R^*\subseteq Q\)。因此，若
\((s,t)\in R^*\) 且 \(P(s)\) 成立，就有 \(P(t)\) 成立；也就是说，
\(P\) 被 \(R^*\) 保持。

\taplsolutionbacklink{ex:2.2.8}{返回练习 2.2.8}
\end{taplsolutionentrybox}
